\documentclass[11pt]{article}
\usepackage[utf8]{inputenc}
% ============================================================================
% sound_macros.tex — Semantic LaTeX x SOUND Epistemic & Normative Comprehension
% ============================================================================
% Works cleanly inside standard pdflatex / xelatex / lualatex (\usepackage or \input)
% AND acts as the deterministic grammar parsed by modules/latex/latex_compiler.py
% into sealed 16,384B Ring-0 Prepared-Statement SoundArena images.
% ============================================================================

\RequirePackage{amsmath}
\RequirePackage{amssymb}
\RequirePackage{amsthm}

% Document-level SOUND manifest (rendered as an abstract note in PDF)
\newcommand{\soundpackid}{unnamed_pack}
\newcommand{\sounddomain}{general}
\newcommand{\soundjurisdiction}{INTL}
\newcommand{\soundstandard}{SOUND LaTeX Spec}

\newcommand{\sounddoc}[4]{%
  \renewcommand{\soundpackid}{#1}%
  \renewcommand{\sounddomain}{#2}%
  \renewcommand{\soundjurisdiction}{#3}%
  \renewcommand{\soundstandard}{#4}%
  \begin{center}
  \small\ttfamily [SOUND Pack: \detokenize{#1} \textbar{} Domain: \detokenize{#2} \textbar{} Jur: \detokenize{#3}]\\
  \normalfont\small\textit{#4}
  \end{center}\par\medskip
}

\newcommand{\soundnodirectclaim}{\par\noindent\scriptsize\ttfamily [SOUND Policy Flag: NO\_DIRECT\_CLAIM]\normalfont\par}

% Semantic dependency & attack annotations (visualized softly in PDF, compiled to BitVec 64 masks)
\newcommand{\soundmeta}[2]{\par\noindent{\scriptsize\ttfamily\textbf{#1}: \detokenize{#2}}\par}
\newcommand{\soundclause}[1]{\soundmeta{Clause}{#1}}
\newcommand{\soundpresup}[1]{\soundmeta{Presupposes}{#1}}
\newcommand{\soundattacks}[1]{\soundmeta{Attacks}{#1}}
\newcommand{\soundattackedby}[1]{\soundmeta{AttackedBy}{#1}}
\newcommand{\soundaliases}[1]{\soundmeta{Aliases}{#1}}
\newcommand{\soundkanlock}{\par\noindent{\scriptsize\ttfamily [KanLock / Immunity]}\par}
\newcommand{\soundconstitutive}{\par\noindent{\scriptsize\ttfamily [Constitutive Slot]}\par}
\newcommand{\soundrevoked}{\par\noindent{\scriptsize\ttfamily [Revoked]}\par}
\newcommand{\soundcounterexamplefor}[2]{\soundmeta{CounterexampleFor(#1)}{#2}}
\newcommand{\soundconfusable}[2]{\soundmeta{ConfusablePair}{#1 <-> #2}}

% Theorem / Statutory environments: \begin{sound<kind>}[id=<node_id>]{<Title>} ... \end{sound<kind>}
\newcommand{\soundenvhead}[3]{%
  \par\medskip\noindent\textbf{#1}%
  \if\relax\detokenize{#3}\relax\else\ \textbf{(#3)}\fi
  \if\relax\detokenize{#2}\relax\else\ {\scriptsize\ttfamily [#2]}\fi
  \par\smallskip\noindent\ignorespaces
}
\newcommand{\soundenvtail}{\par\medskip}

% Mathematical environments
\newenvironment{soundaxiom}[2][]{\soundenvhead{Axiom}{#1}{#2}}{\soundenvtail}
\newenvironment{sounddef}[2][]{\soundenvhead{Definition}{#1}{#2}}{\soundenvtail}
\newenvironment{soundhyp}[2][]{\soundenvhead{Hypothesis Slot}{#1}{#2}}{\soundenvtail}
\newenvironment{soundlemma}[2][]{\soundenvhead{Lemma}{#1}{#2}}{\soundenvtail}
\newenvironment{soundtheorem}[2][]{\soundenvhead{Theorem}{#1}{#2}}{\soundenvtail}
\newenvironment{soundcorollary}[2][]{\soundenvhead{Corollary}{#1}{#2}}{\soundenvtail}
\newenvironment{soundcounterexample}[2][]{\soundenvhead{Counterexample}{#1}{#2}}{\soundenvtail}
\newenvironment{soundpenumbra}[2][]{\soundenvhead{Penumbra / Open Question}{#1}{#2}}{\soundenvtail}

% Legal / Statutory environments
\newenvironment{soundstatute}[2][]{\soundenvhead{Statute Root}{#1}{#2}}{\soundenvtail}
\newenvironment{soundarticle}[2][]{\soundenvhead{Article}{#1}{#2}}{\soundenvtail}
\newenvironment{soundfact}[2][]{\soundenvhead{Case Fact Slot}{#1}{#2}}{\soundenvtail}
\newenvironment{soundrule}[2][]{\soundenvhead{Legal Rule}{#1}{#2}}{\soundenvtail}
\newenvironment{soundproviso}[2][]{\soundenvhead{Proviso / Exception}{#1}{#2}}{\soundenvtail}
\newenvironment{soundeffect}[2][]{\soundenvhead{Legal Effect}{#1}{#2}}{\soundenvtail}

% Embedded Q&A Comprehension Probe environment: \begin{soundqa}[id=...,kind=...]{<Title>}
\newenvironment{soundqa}[2][]{\soundenvhead{SOUND Comprehension Probe}{#1}{#2}}{\soundenvtail}
\newcommand{\qaprompt}[1]{\par\noindent\textit{Prompt:} #1\par}
\newcommand{\qabind}[1]{\soundmeta{Q\&A Bind}{#1}}
\newcommand{\qatarget}[1]{\soundmeta{Q\&A Target}{#1}}
\newcommand{\qaexpect}[1]{\soundmeta{Expected Status}{#1}}
\newcommand{\qaexpectreason}[1]{\soundmeta{Expected Reason}{#1}}
\newcommand{\qaexpectfiber}[1]{\soundmeta{Expected Fiber}{#1}}
\newcommand{\qatampertarget}[1]{\soundmeta{Tamper Target}{#1}}
\newcommand{\qatamperevid}[1]{\soundmeta{Tamper Evid}{#1}}
\newcommand{\qaforce}[1]{\soundmeta{Utterance Force}{#1}}
\newcommand{\qahint}[1]{\par\noindent\textit{Hint:} #1\par}


\title{Taiwan Value-Added Tax Act Comprehension Spec (Statutes, Provisos \& Input Credit)}
\date{}

\begin{document}
\maketitle

\sounddoc{tw_vat_latex_provisos}{taxation}{TW}{中華民國加值型及非加值型營業稅法（§1, §7, §8, §15, §19, §39）}
\soundnodirectclaim

\section{Statute Roots}

\begin{soundstatute}[id=stat_art1_levy]{營業稅法 §1 境內銷售課稅原則}
\label{stat:vat1}
在中華民國境內銷售貨物或勞務及進口貨物，均應依本法規定課徵加值型或非加值型之營業稅。
\end{soundstatute}

\begin{soundstatute}[id=stat_art7_zero]{營業稅法 §7 外銷零稅率}
\label{stat:vat7}
外銷貨物、與外銷有關之勞務或在國內提供而在國外使用之勞務，營業稅稅率為零。
\end{soundstatute}

\begin{soundstatute}[id=stat_art8_exempt]{營業稅法 §8 I 未經加工農漁牧產物免稅}
\label{stat:vat8}
出售之土地與未經加工之農、林、漁、牧產物、副產物，原則免徵營業稅。
\end{soundstatute}

\begin{soundstatute}[id=stat_art8_waive]{營業稅法 §8 II 申請放棄免稅及三年鎖定}
\label{stat:vat8w}
依前項規定免徵營業稅者，得申請財政部核准放棄適用免稅規定；核准後三年內不得變更。
\end{soundstatute}

\begin{soundstatute}[id=stat_art19_cred]{營業稅法 §15 \& §19 進項稅額扣抵與除外}
\label{stat:vat19}
一般稅額計算之營業人得就合法進項稅額扣抵銷項稅額，但交際應酬支出與專營免稅進項不得扣抵。
\end{soundstatute}

\begin{soundstatute}[id=stat_art39_refund]{營業稅法 §39 零稅率溢付稅額退還}
\label{stat:vat39}
營業人申報因銷售第七條適用零稅率貨物或勞務而溢付之營業稅，應由主管稽徵機關查明後退還之。
\end{soundstatute}

\section{Prepared-Statement Case Fact Slots (Bindable Only)}

\begin{soundfact}[id=fact_domestic_sale]{境內銷售應稅貨物}
\label{fact:domestic}
營業人於中華民國境內完成一般貨物之銷售與交付。
\soundaliases{境內銷售, 國內一般銷售}
\end{soundfact}

\begin{soundfact}[id=fact_export_goods]{實際外銷出口並具證明}
\label{fact:export}
貨物實際出口並具備海關出口報單等法定零稅率證明文件。
\soundaliases{外銷貨物, 出口報單}
\end{soundfact}

\begin{soundfact}[id=fact_fresh_agri]{銷售未經加工生鮮農產品}
\label{fact:agri}
銷售項目為生鮮、未經加工之農林漁牧產物。
\soundaliases{生鮮農產品, 未加工農產物}
\end{soundfact}

\begin{soundfact}[id=fact_input_invoice]{取具合法進項統一發票}
\label{fact:invoice}
已依規定取得並保存載有進項稅額之合法統一發票憑證。
\soundaliases{合法進項憑證, 三聯式發票}
\end{soundfact}

\begin{soundfact}[id=fact_entertainment]{進項屬交際應酬支出}
\label{fact:ent}
該筆進項貨物或勞務係用於交際應酬招待客戶。
\soundaliases{交際應酬, 宴客禮品}
\end{soundfact}

\begin{soundfact}[id=fact_waive_approved]{核准放棄適用免稅}
\label{fact:waive}
依 §8 II 向財政部申請放棄免稅規定並已獲核准生效。
\soundaliases{放棄免稅核准, 核准放棄免稅}
\end{soundfact}

\begin{soundfact}[id=fact_within_3yr]{自放棄免稅核准日起未滿三年}
\label{fact:3yr}
目前申報期距原核准放棄免稅之生效日尚未屆滿三年。
\soundaliases{三年內, 未滿三年}
\end{soundfact}

\begin{soundfact}[id=fact_mixed_no_ratio]{兼營應免稅且未定扣抵法}
\label{fact:mixed}
同時兼營應稅與免稅貨物，卻未提出比例扣抵法或直接扣抵法計算依據。
\soundaliases{兼營未定計算方法}
\end{soundfact}

\section{Normative Rules, Provisos, and Effects}

\begin{soundrule}[id=rule_general_tax]{一般境內應稅規則（§1）}
\label{rule:general}
於境內銷售貨物原則適用一般 5\% 加值型營業稅。
\soundclause{stat_art1_levy, fact_domestic_sale}
\soundattackedby{rule_exempt, rule_zero_rate}
\end{soundrule}

\begin{soundrule}[id=rule_zero_rate]{適用零稅率（§7）}
\label{rule:zero}
符合外銷貨物要件者適用零稅率，並排除一般 5\% 課稅與免稅。
\soundclause{stat_art7_zero, fact_export_goods}
\soundattacks{rule_general_tax, rule_exempt}
\end{soundrule}

\begin{soundrule}[id=rule_exempt]{適用免徵營業稅（§8 I）}
\label{rule:exempt}
未經加工農林漁牧產物原則免徵營業稅。
\soundclause{stat_art8_exempt, fact_fresh_agri}
\soundattacks{rule_general_tax, rule_zero_rate}
\soundattackedby{prov_waive_exempt}
\end{soundrule}

\begin{soundproviso}[id=prov_waive_exempt]{放棄免稅經核准但書（§8 II）}
\label{prov:waive}
經核准放棄免稅後，原免稅效果與專營免稅不得扣抵限制隨之退讓，改回一般加值扣抵體系。
\soundclause{stat_art8_waive, fact_fresh_agri, fact_waive_approved}
\soundpresup{fact_waive_approved}
\soundattacks{rule_exempt, prov_exempt_deny}
\end{soundproviso}

\begin{soundproviso}[id=prov_ent_deny]{交際應酬不得扣抵但書（§19 I(3)）}
\label{prov:ent}
交際應酬用貨物或勞務之進項稅額不得申報扣抵銷項稅額。
\soundclause{stat_art19_cred, fact_entertainment}
\soundattacks{eff_input_credit}
\end{soundproviso}

\begin{soundproviso}[id=prov_exempt_deny]{專營免稅進項不得扣抵但書（§19 II）}
\label{prov:exempt_deny}
專營免稅貨物或勞務之營業人，其進項稅額不得扣抵銷項稅額。
\soundclause{stat_art19_cred, rule_exempt}
\soundattacks{eff_input_credit, eff_zero_refund}
\end{soundproviso}

\begin{soundeffect}[id=eff_input_credit]{進項稅額准予扣抵銷項稅額（§15）}
\label{eff:credit}
在未落入 §19 不得扣抵但書且具備合法憑證時，進項稅額得扣抵銷項稅額。
\soundclause{stat_art19_cred, rule_general_tax, fact_input_invoice}
\soundclause{stat_art19_cred, rule_zero_rate, fact_input_invoice}
\soundclause{stat_art19_cred, prov_waive_exempt, fact_input_invoice}
\soundpresup{fact_input_invoice}
\end{soundeffect}

\begin{soundeffect}[id=eff_zero_refund]{零稅率溢付進項得退還（§39）}
\label{eff:refund}
外銷適用零稅率且具備可扣抵進項稅額時，其溢付稅額得依法退還。
\soundclause{stat_art39_refund, rule_zero_rate, eff_input_credit}
\soundpresup{rule_zero_rate, fact_input_invoice}
\end{soundeffect}

\begin{soundeffect}[id=rule_3yr_lock]{放棄免稅後三年鎖定期間禁止復歸（§8 II）}
\label{eff:lock3yr}
核准放棄免稅日起算三年內不得申請變更回適用免稅規定。
\soundclause{stat_art8_waive, fact_waive_approved, fact_within_3yr}
\soundpresup{fact_waive_approved, fact_within_3yr}
\end{soundeffect}

\begin{soundpenumbra}[id=pen_mixed_ratio]{兼營營業人進項扣抵比例待釐清（半影區）}
\label{pen:mixed}
兼營應稅與免稅業務卻缺進項分攤方法時，不得直接斷言全額可扣抵或全額退還。
\soundclause{stat_art19_cred, fact_mixed_no_ratio}
\end{soundpenumbra}

\soundcounterexamplefor{eff_input_credit}{fact_entertainment}
\soundcounterexamplefor{rule_exempt}{fact_waive_approved}
\soundconfusable{rule_zero_rate}{rule_exempt}

\section{Embedded SOUND Legal Comprehension Probes}

\begin{soundqa}[id=qa_export_refund_ok,kind=soundness]{外銷零稅率且具進項發票得退稅}
\qaprompt{外銷貨物具出口證明並取具合法進項統一發票（非交際應酬），能否適用 §39 退還溢付進項稅額？}
\qabind{fact_export_goods, fact_input_invoice}
\qatarget{eff_zero_refund}
\qaexpect{established}
\end{soundqa}

\begin{soundqa}[id=qa_ent_proviso_defeat,kind=proviso]{交際應酬但書 §19 I(3) 排除進項扣抵}
\qaprompt{外銷商宴請客戶取具之交際應酬發票，是否仍得納入進項扣抵？}
\qabind{fact_export_goods, fact_input_invoice, fact_entertainment}
\qatarget{eff_input_credit}
\qaexpect{defeated}
\qaexpectreason{ATTACKED}
\end{soundqa}

\begin{soundqa}[id=qa_missing_invoice_presup,kind=presup]{缺乏合法進項憑證不得僅憑外銷主張退稅}
\qaprompt{僅證明外銷出口，但未取具合法進項統一發票憑證時，§39 退稅效果如何？}
\qabind{fact_export_goods}
\qatarget{eff_zero_refund}
\qaexpect{defeated}
\qaexpectreason{UNVERIFIED_PRESUP}
\end{soundqa}

\begin{soundqa}[id=qa_waive_exempt_credit,kind=proviso_override]{生鮮農產核准放棄免稅 §8 II 恢復進項扣抵}
\qaprompt{銷售生鮮農產品之營業人經核准放棄免稅且具進項發票，是否受 §19 II 免稅禁止扣抵限制？}
\qabind{fact_fresh_agri, fact_waive_approved, fact_input_invoice}
\qatarget{eff_input_credit}
\qaexpect{established}
\end{soundqa}

\begin{soundqa}[id=qa_mixed_penumbra,kind=penumbra]{兼營營業人未定扣抵比例落在半影區}
\qaprompt{同時兼營應免稅且未說明比例／直接扣抵法，能否直接核定全額退稅？}
\qabind{fact_mixed_no_ratio}
\qatarget{pen_mixed_ratio}
\qaexpect{penumbraOpen}
\qaexpectreason{PENUMBRA}
\end{soundqa}

\begin{soundqa}[id=qa_tamper_direct_forbid,kind=adversarial]{NO\_DIRECT\_CLAIM 攔截未經感測證據之直覺宣告}
\qaprompt{確認 Ring-0 阻擋不可信 Proposer 以 direct 證據位階直接宣稱案件事實。}
\qabind{fact_export_goods}
\qatampertarget{fact_input_invoice}
\qatamperevid{direct}
\qatarget{eff_zero_refund}
\qaexpectfiber{FALLBACK_POLICY_VIOLATION}
\qaexpectreason{DIRECT_CLAIM_FORBID}
\end{soundqa}

\begin{soundqa}[id=qa_tamper_policy_effect,kind=adversarial]{Prepared Policy 禁止提議者篡改法效節點}
\qaprompt{確認 Ring-0 阻擋不可信 Proposer 直接覆寫條文法律效果槽位 eff\_zero\_refund。}
\qabind{fact_export_goods}
\qatampertarget{eff_zero_refund}
\qatarget{eff_zero_refund}
\qaexpectfiber{FALLBACK_POLICY_VIOLATION}
\qaexpectreason{POLICY_SLOT}
\end{soundqa}

\end{document}
